\section{Method}

Figure~\ref{fig:overview} shows the setup. Four knowledge sources feed one
slot in the prompt of an otherwise fixed refinement loop; the loop's
generators are scored by a differential harness over four deserialization
paths.

\begin{figure}[t]
\centering
\resizebox{\columnwidth}{!}{%
\begin{tikzpicture}[
  font=\footnotesize,
  box/.style={draw, rounded corners=2pt, align=center, inner sep=3pt},
  src/.style={box, minimum width=1.55cm, minimum height=0.9cm},
  arr/.style={-{Stealth[length=4pt]}, thick}]
\node[src] (none)  at (0,3.0)  {\textbf{none}\\empty};
\node[src] (human) at (0,1.95) {\textbf{human}\\4 observations};
\node[src] (code)  at (0,0.9)  {\textbf{code}\\agent reads\\source};
\node[src] (probe) at (0,-0.25) {\textbf{probe}\\agent runs\\20 probes};
\node[box, minimum width=1.5cm, minimum height=1.1cm] (slot) at (2.55,1.35) {knowledge\\slot\\$\le$1{,}500 chars};
\node[box, minimum width=2.1cm, minimum height=1.9cm] (loop) at (4.95,1.35)
  {\textbf{refinement loop}\\(fixed)\\[2pt]LLM writes a\\Hypothesis generator\\$\to$ 500 documents\\$\times$ 5 iterations};
\node[box, minimum width=1.9cm, minimum height=1.9cm] (harness) at (8.05,1.35)
  {\textbf{harness}\\4 paths\\[2pt]divergence\\oracle};
\foreach \n in {none,human,code,probe} \draw[arr] (\n.east) -- (slot.west);
\draw[arr] (slot) -- (loop);
\draw[arr] ([yshift=4pt]loop.east) -- node[above, font=\scriptsize]{docs} ([yshift=4pt]harness.west);
\draw[arr] ([yshift=-4pt]harness.west) -- node[below, font=\scriptsize]{score} ([yshift=-4pt]loop.east);
\draw[arr, dashed] (probe.south) |- ++(0,-0.35) -| node[pos=0.25, below, font=\scriptsize]{probe documents and per-path outcomes} (harness.south);
\end{tikzpicture}}
\caption{Study design. Only the knowledge slot differs between arms. The
probing agent runs its probes through the same harness the loop is scored
by.}
\label{fig:overview}
\end{figure}

\subsection{Targets and oracle}

Both targets decode the same logical record: a non-null integer
\texttt{id}, a string \texttt{amount}, a nullable string \texttt{name}, a
\texttt{status} enum with three values, a non-null list of strings
\texttt{tags}, and a nullable recursive \texttt{child}. Each generated
document is run through four deserialization paths in one
long-lived harness process. A document \emph{diverges} if one path accepts
it and another rejects it, or if all accept but decode different values.
Values are compared through a canonical JSON rendering of the decoded
object built only from primitives, so that the comparison itself cannot
fail. The harness speaks newline-delimited JSON on standard input and output,
with each document base64-encoded, so arbitrary bytes reach the paths
unchanged. The primary metric, as in the previous study, is the percentage
of all generated documents that diverge.

\textbf{Dart} reuses the previous study's harness unchanged~\dartcite: hand-written
parsing, \texttt{json\_serializable}~\cite{jsonserializable2026},
\texttt{freezed}~\cite{freezed2026} and
\texttt{built\_value}~\cite{builtvalue2026}, on Dart
3.13.2~\cite{dart2026lang}. All four receive the output of one shared
\texttt{jsonDecode}, so invalid JSON is rejected identically before any
path runs. Two divergence patterns are known (written in path order):
\texttt{RAAR}, where the two paths that decode \texttt{id} through
\texttt{num.toInt()} accept a large integer literal that
\texttt{jsonDecode} has turned into a double, and saturate it; and
\texttt{RRRA}, where \texttt{built\_value} alone accepts a document with
\texttt{tags} missing and substitutes an empty list.

\textbf{Kotlin} uses a new harness with the same protocol and four libraries
at their documented defaults: Gson~\cite{gson2026},
Moshi with its generated adapter~\cite{moshi2026},
kotlinx.serialization~\cite{kotlinxserialization2026} and Jackson with its
Kotlin module~\cite{jacksonkotlin2026}. The record is an idiomatic data
class with no default values and no configuration beyond the annotations each
library needs to run. Unlike Dart, each library parses the raw text itself,
so all four run on every input and syntax-level divergence counts. The known
answer key is the Gson null-safety bypass~\cite{gsonissue1657}: Gson
constructs objects without Kotlin's constructor checks, so a missing or
\texttt{null} value for a non-null property is accepted as \texttt{null}.
Divergence that follows from a documented default is labeled as such and
never called a bug.

\subsection{The refinement loop}

The loop is the previous study's score-only loop~\dartcite, unchanged
except for one setting. Each of five iterations asks \texttt{gpt-4.1-mini}~\cite{openai2025gpt41}
(temperature 0.2) for a Hypothesis generator of JSON documents, draws 500
documents from it in a restricted sandbox, and reports back how many
diverged. The prompt describes the schema and the four paths, gives the
score and a generic strategy hint (vary one thing at a time: a wrong type, a
missing field, a boundary value), and asks for a new generator. A test
rebuilds all 100 prompts committed by the previous study's score-only and
knowledge arms byte for byte. The one change is the output-token limit,
raised from 2{,}500 to 8{,}000 for every arm before any experiment data, after
a pilot lost four of five generators to truncation; no response in the main
study was truncated, and one of the extension's 75 was.

\subsection{Knowledge arms}

Each arm fills one knowledge slot in the prompt: a one-sentence header naming
the source, followed by a body of at most 1{,}500 characters.

\textbf{None.} The slot is empty.

\textbf{Human} (Dart only). The four observations from the previous study,
verbatim (1{,}291 characters): that invalid JSON never reaches the paths;
that \texttt{built\_value} alone accepts a missing \texttt{tags} and uses
an empty list; that wrong types and nulls are rejected by all four in every
case tried; and that two paths decode \texttt{id} as \texttt{num.toInt()},
which saturates for integer literals outside the 64-bit range. The label
means that the researchers produced this knowledge outside the pipeline,
from documents they chose, rather than an agent inside it.

\textbf{Code} (Dart only). An agent reads a fixed, per-seed identical set
of nine source files, about 17{,}000 characters: each path's model and its
generated \texttt{*.g.dart} code. It then writes a numbered list of
behavioral observations. For Kotlin no comparable source exists: Gson and
Jackson are reflective and kotlinx.serialization generates bytecode with a
compiler plugin. That asymmetry is itself a practical argument for black-box
probing.

\textbf{Probe.} An agent designs up to 20 probe documents in two rounds, at
most 12 then the rest, written as raw text so that non-JSON probes are
possible. Each probe runs through the real harness, and the agent sees, per
path, whether it accepted or rejected the document, the exception type, and
the decoded value. After the second round it writes a numbered list of
observations.

Both agents use the same model and settings as the loop and see the same
target description and strategy hint as the \emph{none} arm, so each starts
from exactly the \emph{none} arm's a priori information. If a summary is
over the cap, one shortening call follows, then a cut at the last line break
under the cap; the cut applied in 19 of 20 acquisitions (mean body 1{,}422
characters). Knowledge is acquired once per seed, so its quality varies
across seeds, and that variance is part of the result.

\subsection{Extension: two remedies for probing}
\label{sec:extension}

After the main results, we declared an extension in the deviations log and
published it before any of its runs, with two new arms and no new seeds for
any existing arm.

\textbf{Systematic probing} (\emph{sysprobe}; Dart and Kotlin). The probing
agent, model, budget and cap are unchanged. Its prompt adds the checklist in
Fig.~\ref{fig:checklist}, and its summary must end with a line listing the
categories no probe covered. Every item is a standard, schema-independent
category from category-partition testing~\cite{ostrand1988category} and
boundary-value analysis~\cite{myers2011art}, and
none names a field, path or library; we disclose that it was written by an
author who knew the answer key. The knowledge header is the same as for
\emph{probe}, so only the body differs.

\textbf{A larger prober} (\emph{probe-gpt41}; Dart). The unchanged probing
agent and prompts, with \texttt{gpt-4.1} for knowledge acquisition only; the
loop stays \texttt{gpt-4.1-mini}.

Both arms use the seeds of the corresponding \emph{probe} arm. The declared
tests are systematic probing against none on Dart (primary), against unguided
probing, human and code on Dart, and against none and unguided probing on
Kotlin, with a Holm correction across these six; and the larger prober
against unguided probing and against none, with a Holm correction across
the two.

\begin{figure}[t]
\centering
\fbox{\begin{minipage}{0.94\columnwidth}
\footnotesize
Probe systematically. Spend your probes on these categories first; they apply
to any JSON schema, and each probe should change one thing in an otherwise
valid document:
(1) a required field missing;
(2) a non-nullable field set to null;
(3) a field with a value of the wrong JSON type;
(4) a numeric field at a representation boundary: a fraction, an exponent
form, and integers just beyond the 32-, 53- and 64-bit ranges;
(5) an enum-like field with an unknown value, and with a case variant;
(6) an extra unknown key, and a duplicate key;
(7) the same kinds of change inside a nested object;
(8) text that is not strict JSON: single quotes, an unquoted value, a
trailing comma, NaN, and a raw control character inside a string.
Prefer breadth across categories over repetition.
\end{minipage}}
\caption{The checklist added to the systematic probing agent's prompt
(abridged; the full text is in the artifact).}
\label{fig:checklist}
\end{figure}

\subsection{Design and analysis}

The study was pre-registered in a design document committed before any code
or API spend; later changes are listed in a dated deviations log. Dart arms
run on seeds 700--704 and Kotlin arms on 800--804. Pilots used seeds outside
these ranges, and the probing and code agents' prompts were frozen after one
pilot change: both agents were given the refinement prompt's existing
strategy hint, which the \emph{none} arm already sees.

For each run, iterations whose generator failed validation contribute no
documents, as in the previous study. We compare arms with an exact two-sided
Mann--Whitney U test and report the mean difference and Cliff's $\delta$.
The primary test is \emph{probe} against \emph{none} on Dart, with
\emph{probe} against \emph{human} co-primary; all six Dart pairs are reported
with a Holm correction alongside raw $p$-values. The Kotlin primary is
\emph{probe} against \emph{none}. A $p \geq 0.05$ is reported as ``not
detected at $n{=}5$'', not as ``no effect''. A claim of comparability to
human knowledge requires $|\delta| < 0.33$ in addition to $p \geq 0.05$.
Replication on fresh seeds was allowed only if declared in the log before
it ran. Secondary measures are pattern recall (the fraction of runs that
find each known pattern at least once) and cost in calls and tokens.
